\subsection{Image-dataset validation}
\label{methods:image_validation}

We evaluated MNIST, Fashion-MNIST and CIFAR-10 using their official training
and test partitions \cite{lecun1998gradient,xiao2017fashion,krizhevsky2009learning}.
Dataset archives were obtained from the PyTorch MNIST mirror, the official
Fashion-MNIST repository and the University of Toronto CIFAR-10 server,
respectively. Archive checksums and download locations are provided with the
code. From each official training partition, a stratified sample of 5,000
images was withheld using seed 410001. Its randomly ordered indices were
partitioned into 4,000 images for model validation, 512 for operator
calibration and 488 for calibration validation. The remaining 55,000 images
for each grayscale dataset and 45,000 CIFAR-10 images were used for model
training. Channel means and standard deviations were computed from these
model-training images only. Official test images were excluded from model
training, checkpoint selection and operator calibration.

Each image was divided into a $4\times4$ grid of non-overlapping patches,
yielding 16 tokens. Patch sizes were $7\times7$ for the $28\times28$
grayscale images and $8\times8$ for the $32\times32$ RGB images; no image
downsampling was applied. A learned linear map projected each flattened
patch to 32 features, followed by a learned positional embedding. All
datasets used four pre-normalization Transformer blocks with two attention
heads, a head dimension of 16 and a feed-forward width of 64. Each block
contained both residual connections and a GELU activation. A final layer
normalization, mean pooling across tokens and a ten-class linear head
produced image-level logits. The reference networks contained 36,682
parameters for each grayscale dataset and 41,258 parameters for CIFAR-10.
Each model-specific HI-SMURF operator added 16,512 parameters, shared across
all blocks and heads.

Three reference models were trained from scratch for each dataset using
initialization seeds 11, 22 and 33. AdamW used a batch size of 128, an initial
learning rate of $10^{-3}$, weight decay of 0.01 and gradient clipping at
1. The learning-rate schedule used a 200-step warm-up followed by cosine
decay to 5\% of the initial rate. Training comprised 6,000 steps for MNIST,
8,000 for Fashion-MNIST and 12,000 for CIFAR-10. Augmentation consisted of
random translations of up to two pixels for grayscale images and four
pixels for CIFAR-10, with zero and reflection padding, respectively.
Fashion-MNIST and CIFAR-10 additionally used random horizontal flips.
CIFAR-10 training used label smoothing of 0.05. Checkpoints were selected
by model-validation accuracy, using cross-entropy to resolve ties. No
pretraining or task-level fine-tuning was applied during conversion to FROST.

For each frozen reference model, attention score rows from the 512
calibration images were collected across all four blocks and both heads.
One shared categorical HI-SMURF operator, with 16 states per input and
TT rank 8, was initialized from the preceding softmax study. We optimized
its core logits for 1,200 Adam steps at each encoding bound in
$\{2,4,8\}$, using batches of 512 score rows. The loss was 256 times the
mean squared error against the reference softmax probabilities. The
learning rate decayed from 0.01 to 0.001, with gradient clipping at 5.
The bound and checkpoint minimizing probability RMSE on the separate
calibration-validation images were retained. Calibration did not use task
labels, and the Transformer parameters remained fixed. Each dataset and
trained-model seed had its own calibrated operator; this experiment did
not evaluate zero-shot transfer of one calibrated operator across datasets.

Finite-stream inference used actual sign--magnitude products throughout
the learned linear stages and attention products. Packed AND and population
counts accelerated simulation while preserving the signed counts of the
explicit bit products exactly. Softmax inputs drove independent saturating
FSMs with 256 burn-in cycles. Because the bounded categorical TT cores are
row stochastic, their conditional output law was sampled by drawing the
successive latent bond indices and then the output category. This ancestral
procedure avoids explicitly contracting a full output probability vector
for every cycle; its distribution was checked against direct contraction.
GELU and reciprocal-square-root gates used eight and four states,
respectively, with 64 burn-in cycles. Layer normalization used stochastic
variance estimation. Arithmetic streams used independently digitally
shifted coordinates of a scrambled two-dimensional Sobol sequence; FSM
drives remained IID Bernoulli. Inference batches contained at most 16
images, sharing weight streams within each batch. Input normalization,
centering, pooling, residual and bias addition, static normalization-affine
gains, sorting, range scaling and stage conversions remained numerical.
The reported stream length excludes burn-in and does not represent
whole-network latency.

The main comparison evaluated all 10,000 official test images per model
using floating-point inference, a simulated W8A8 baseline, analytic FROST
and finite-stream FROST at $L=256$. The finite-stream evaluation used the
prespecified stream seed 101, with deterministic offsets for batches. W8A8
used dynamic per-token activation and per-output-channel weight symmetric
int8 quantization for learned linear maps and query--key products, with
int32 accumulation. Attention probabilities used unsigned eight-bit
quantization and values used per-channel symmetric int8 quantization for
attention aggregation. Softmax, GELU, layer normalization and residual
operations remained floating point in this baseline. This was a numerical
precision control, with no hardware speed, area or energy comparison.

A separate diagnostic evaluation used a fixed, randomly selected subset
of 1,000 test images, comprising 100 images per class (seed 420001).
This subset compared $L\in\{64,256,1024\}$ using stream seed 101 and
direct transfer of the uncalibrated HI-SMURF operator at $L=256$.
The calibrated $L=256$, seed-101 predictions were taken from the full-test
evaluation. Additional subset evaluations at seeds 202 and 303 assessed
stream variability. Main results and stream-length plots report means and
sample standard deviations across three trained models at seed 101.
For the separate stream-repeat summary, results were first averaged over
three stream realizations within each model before computing between-model
means and sample standard deviations. These subset results were kept
separate from full-test estimates.

We reported classification accuracy, macro-averaged F1, negative
log-likelihood, agreement with the floating-point prediction, paired
accuracy differences in percentage points and relative centred-logit
error. For reference and approximate logit matrices $Z$ and $\widehat Z$,
let $C(Z)$ subtract the mean of each row. The latter error was
$\|C(\widehat Z)-C(Z)\|_{\mathrm F}/\|C(Z)\|_{\mathrm F}$,
evaluated over all images in the indicated split. Per-image predictions,
test indices, model checkpoints, calibration records and figure source
data are included in the accompanying reproducibility files.
